\documentclass[11pt,a4paper]{article}
\usepackage[utf8]{inputenc}
\usepackage[T1]{fontenc}
\usepackage[spanish,es-minimal,es-noshorthands]{babel}
\usepackage{lmodern}
\usepackage[sfdefault,lining]{FiraSans}
\usepackage[a4paper,top=2.3cm,bottom=2.2cm,left=1.8cm,right=1.8cm,headheight=1.3cm,headsep=0.4cm,footskip=1cm]{geometry}
\usepackage{amsmath,amssymb}
\allowdisplaybreaks
\usepackage{icomma}
\usepackage{graphicx}
\usepackage[table]{xcolor}
\usepackage{tikz}
\usepackage{fontawesome5}
\usepackage{ccicons}
\usepackage{fancyhdr}
\usepackage{lastpage}
\usepackage{tasks}
\usepackage{needspace}
\usepackage[most]{tcolorbox}
\usepackage{array}
\usepackage[hidelinks]{hyperref}
\AtBeginDocument{\hypersetup{pdftitle={Ficha de ejercicios de la unidad 3: Vectores y aplicaciones. Rectas, planos y haces de planos. 2º Bach.},pdfsubject={Matemáticas II 2º Bach.},pdfkeywords={vectores, producto escalar, producto vectorial, producto mixto, áreas, volúmenes, ecuaciones de la recta, ecuaciones del plano, posiciones relativas, haces de planos, PAU}}}

\definecolor{navy}{HTML}{1B2748}
\definecolor{accent}{HTML}{E8603A}
\definecolor{teal}{HTML}{1F9E90}
\definecolor{soft}{HTML}{EEF1F8}
\definecolor{muted}{HTML}{5D6685}

\setlength{\parindent}{0pt}
\setlength{\parskip}{3pt}
\settasks{label=\alph*), label-format=\bfseries\color{navy}, label-width=1.3em, item-indent=1.8em, column-sep=0.6em, after-item-skip=5pt, before-skip=3pt, after-skip=4pt}

% ---------- capçalera i peu ----------
\newcommand{\llicencia}{\href{https://creativecommons.org/licenses/by-nc-sa/4.0/deed.es}{\ccbyncsa}\ \ Licencia Creative Commons Atribución-NoComercial-CompartirIgual 4.0 Internacional (CC BY-NC-SA 4.0)}
\pagestyle{fancy}
\fancyhf{}
\fancyhead[L]{\raisebox{-0.35\height}{\includegraphics[height=0.95cm]{logo-gva.png}}}
\fancyhead[C]{\small\color{muted}Ficha de ejercicios · Unidad 3: Vectores, rectas y planos · 2\textsuperscript{o} Bach.}
\fancyhead[R]{\raisebox{-0.35\height}{\includegraphics[height=0.95cm]{logo-vermellar.png}}}
\fancyfoot[L]{\scriptsize\color{muted}\llicencia}
\fancyfoot[R]{\small\color{muted}\thepage}
\renewcommand{\headrulewidth}{0.4pt}
\renewcommand{\headrule}{\color{muted!40}\hrule height \headrulewidth}
\fancypagestyle{primera}{\fancyhf{}\renewcommand{\headrulewidth}{0pt}
  \fancyfoot[L]{\scriptsize\color{muted}\llicencia}
  \fancyfoot[R]{\small\color{muted}\thepage}}

% ---------- nivells i exercicis ----------
\newcommand{\estrelles}[1]{{\color{accent}\ifcase#1\or\faStar\faStar[regular]\faStar[regular]\or\faStar\faStar\faStar[regular]\or\faStar\faStar\faStar\fi}}
\newcounter{exer}
\newcommand{\ex}[1]{\par\needspace{4\baselineskip}\medskip\refstepcounter{exer}\noindent{\color{navy}\bfseries\theexer.}\ \ }

\newcommand{\seccio}[2]{\par\needspace{10\baselineskip}\bigskip
  \begin{tcolorbox}[colback=navy,colframe=navy,coltext=white,arc=2mm,boxrule=0pt,left=3mm,right=3mm,top=1.2mm,bottom=1.2mm]
  {\large\bfseries #1\quad #2}
  \end{tcolorbox}\vspace{-2pt}}


\setlength{\emergencystretch}{3em}  % les fórmules llargues (equacions de rectes, ternes) passen a la línia següent

% quadre de recordatori (només a la fitxa d'exercicis)
\newcommand{\recorda}[1]{\par\smallskip\begin{tcolorbox}[colback=accent!7,colframe=accent!55,arc=2mm,boxrule=0.5pt,left=2.5mm,right=2.5mm,top=1mm,bottom=1mm]\small{\color{accent}\faLightbulb[regular]}\ \textbf{Recuerda:}\ #1\end{tcolorbox}}
% etiqueta dels exercicis de la PAU d'altres comunitats
\newcommand{\pau}[1]{{\color{teal}\small\bfseries\faUniversity\ PAU · #1.}\ }
\newcommand{\res}[1]{#1}\newcommand{\rest}[1]{#1}

% =====================================================================
% Capçalera FentMates de les fitxes: un color per a cada curs.
% Els colors són els de la pàgina inicial (LEVELS a assets/js/materials.js);
% el logo ve de _eines/logo/pdf/ (python3 _eines/logo/genera.py).
%
% Ús, al final del preàmbul de la unitat:
%   % =====================================================================
% Capçalera FentMates de les fitxes: un color per a cada curs.
% Els colors són els de la pàgina inicial (LEVELS a assets/js/materials.js);
% el logo ve de _eines/logo/pdf/ (python3 _eines/logo/genera.py).
%
% Ús, al final del preàmbul de la unitat:
%   % =====================================================================
% Capçalera FentMates de les fitxes: un color per a cada curs.
% Els colors són els de la pàgina inicial (LEVELS a assets/js/materials.js);
% el logo ve de _eines/logo/pdf/ (python3 _eines/logo/genera.py).
%
% Ús, al final del preàmbul de la unitat:
%   \input{../comu/fentmates.tex}
%   \fmcurs{1eso}
% i, al començament del document:
%   \capcalera{Matemàtiques}{1r ESO}{Fitxa d'exercicis · Unitat 2}{Divisibilitat}
% Opcional abans de \input: \def\fmneed{14} (espai mínim abans de cada secció).
% =====================================================================
\makeatletter
\def\fm@logodir{../../logo/pdf}
\providecommand{\fmneed}{10}
\def\fm@cursnom{}

% clau del curs, color i color del text damunt del color
\newcommand{\fm@curs}[3]{\@namedef{fm@color@#1}{#2}\@namedef{fm@text@#1}{#3}}
\fm@curs{1eso}{E6B84A}{navy}
\fm@curs{2eso}{3B82F6}{white}
\fm@curs{3eso}{FF7A50}{white}
\fm@curs{4eso-a}{A78BFA}{white}
\fm@curs{4eso-b}{6D28D9}{white}
\fm@curs{1batx-mi}{9B1C1C}{white}
\fm@curs{1batx-ccss}{EF6B6B}{white}
\fm@curs{2batx-mii}{1E6B47}{white}
\fm@curs{2batx-ccss}{2BB3A3}{white}

\newcommand{\fmcurs}[1]{%
  \@ifundefined{fm@color@#1}{\PackageError{fentmates}{Curs desconegut: #1}{}}{}%
  \edef\fm@tmp{\noexpand\definecolor{curs}{HTML}{\@nameuse{fm@color@#1}}%
               \noexpand\colorlet{curstext}{\@nameuse{fm@text@#1}}}\fm@tmp
  \colorlet{cursfosc}{curs!72!black}%
  \def\fm@logo{\fm@logodir/vertical-#1.pdf}%
  \def\fm@marca{\fm@logodir/marca-#1.pdf}}

% ---------- capçalera de la primera pàgina ----------
% #1 matèria (petit), #2 curs, #3 tipus de document i unitat, #4 títol
\newcommand{\capcalera}[4]{%
  \gdef\fm@cursnom{#2}%
  \par\noindent
  \begin{tcolorbox}[colback=curs,colframe=curs,boxrule=0pt,arc=2.2mm,width=3.5cm,height=2.15cm,
      valign=center,left=3mm,right=2mm,top=1mm,bottom=1mm,nobeforeafter,box align=center]
    \color{curstext}{\footnotesize #1}\par\vspace{1pt}{\fontsize{19}{21}\selectfont\bfseries #2\par}%
  \end{tcolorbox}\hspace{5mm}%
  \begin{minipage}[c]{\dimexpr\linewidth-3.5cm-5mm-2.9cm\relax}
    {\color{cursfosc}\small\bfseries\MakeUppercase{#3}}\par\vspace{2pt}
    {\color{navy}\fontsize{20}{23}\selectfont\bfseries #4\par}
  \end{minipage}\hfill
  \begin{minipage}[c]{2.6cm}\raggedleft\includegraphics[width=2.6cm]{\fm@logo}\end{minipage}\par
  \vspace{2pt}}

% ---------- barres de secció del color del curs ----------
\renewcommand{\seccio}[2]{\par\needspace{\fmneed\baselineskip}\bigskip
  \begin{tcolorbox}[colback=curs,colframe=curs,coltext=curstext,arc=2mm,boxrule=0pt,left=3mm,right=3mm,top=1.2mm,bottom=1.2mm]
  {\large\bfseries #1\quad #2}
  \end{tcolorbox}\vspace{-2pt}}

% ---------- estrelles de dificultat del color del curs (a Batxillerat no se'n posen) ----------
\@ifundefined{estrelles}{}{%
  \renewcommand{\estrelles}[1]{{\color{curs}\ifcase#1\or\faStar\faStar[regular]\faStar[regular]\or\faStar\faStar\faStar[regular]\or\faStar\faStar\faStar\fi}}}

% ---------- capçalera i peu de les pàgines ----------
% primera pàgina: logos institucionals al peu
\fancypagestyle{primera}{\fancyhf{}\renewcommand{\headrulewidth}{0pt}
  \fancyfoot[L]{\raisebox{-0.32\height}{\includegraphics[height=0.78cm]{logo-gva.png}}\hspace{5mm}%
    \raisebox{-0.32\height}{\includegraphics[height=0.78cm]{logo-vermellar.png}}}
  \fancyfoot[R]{\scriptsize\color{muted}\href{https://creativecommons.org/licenses/by-nc-sa/4.0/deed.ca}{\ccbyncsa}\ \ CC BY-NC-SA 4.0 · fentmates.org\qquad\small\thepage}}
% la resta: el símbol a l'esquerra i el curs a la dreta, amb el títol al centre com abans
\fancyhead[L]{\raisebox{-0.3\height}{\includegraphics[height=0.62cm]{\fm@marca}}}
\fancyhead[R]{\small\bfseries\color{cursfosc}\fm@cursnom}
\renewcommand{\headrule}{\color{curs!45}\hrule height \headrulewidth}
\makeatother

%   \fmcurs{1eso}
% i, al començament del document:
%   \capcalera{Matemàtiques}{1r ESO}{Fitxa d'exercicis · Unitat 2}{Divisibilitat}
% Opcional abans de \input: \def\fmneed{14} (espai mínim abans de cada secció).
% =====================================================================
\makeatletter
\def\fm@logodir{../../logo/pdf}
\providecommand{\fmneed}{10}
\def\fm@cursnom{}

% clau del curs, color i color del text damunt del color
\newcommand{\fm@curs}[3]{\@namedef{fm@color@#1}{#2}\@namedef{fm@text@#1}{#3}}
\fm@curs{1eso}{E6B84A}{navy}
\fm@curs{2eso}{3B82F6}{white}
\fm@curs{3eso}{FF7A50}{white}
\fm@curs{4eso-a}{A78BFA}{white}
\fm@curs{4eso-b}{6D28D9}{white}
\fm@curs{1batx-mi}{9B1C1C}{white}
\fm@curs{1batx-ccss}{EF6B6B}{white}
\fm@curs{2batx-mii}{1E6B47}{white}
\fm@curs{2batx-ccss}{2BB3A3}{white}

\newcommand{\fmcurs}[1]{%
  \@ifundefined{fm@color@#1}{\PackageError{fentmates}{Curs desconegut: #1}{}}{}%
  \edef\fm@tmp{\noexpand\definecolor{curs}{HTML}{\@nameuse{fm@color@#1}}%
               \noexpand\colorlet{curstext}{\@nameuse{fm@text@#1}}}\fm@tmp
  \colorlet{cursfosc}{curs!72!black}%
  \def\fm@logo{\fm@logodir/vertical-#1.pdf}%
  \def\fm@marca{\fm@logodir/marca-#1.pdf}}

% ---------- capçalera de la primera pàgina ----------
% #1 matèria (petit), #2 curs, #3 tipus de document i unitat, #4 títol
\newcommand{\capcalera}[4]{%
  \gdef\fm@cursnom{#2}%
  \par\noindent
  \begin{tcolorbox}[colback=curs,colframe=curs,boxrule=0pt,arc=2.2mm,width=3.5cm,height=2.15cm,
      valign=center,left=3mm,right=2mm,top=1mm,bottom=1mm,nobeforeafter,box align=center]
    \color{curstext}{\footnotesize #1}\par\vspace{1pt}{\fontsize{19}{21}\selectfont\bfseries #2\par}%
  \end{tcolorbox}\hspace{5mm}%
  \begin{minipage}[c]{\dimexpr\linewidth-3.5cm-5mm-2.9cm\relax}
    {\color{cursfosc}\small\bfseries\MakeUppercase{#3}}\par\vspace{2pt}
    {\color{navy}\fontsize{20}{23}\selectfont\bfseries #4\par}
  \end{minipage}\hfill
  \begin{minipage}[c]{2.6cm}\raggedleft\includegraphics[width=2.6cm]{\fm@logo}\end{minipage}\par
  \vspace{2pt}}

% ---------- barres de secció del color del curs ----------
\renewcommand{\seccio}[2]{\par\needspace{\fmneed\baselineskip}\bigskip
  \begin{tcolorbox}[colback=curs,colframe=curs,coltext=curstext,arc=2mm,boxrule=0pt,left=3mm,right=3mm,top=1.2mm,bottom=1.2mm]
  {\large\bfseries #1\quad #2}
  \end{tcolorbox}\vspace{-2pt}}

% ---------- estrelles de dificultat del color del curs (a Batxillerat no se'n posen) ----------
\@ifundefined{estrelles}{}{%
  \renewcommand{\estrelles}[1]{{\color{curs}\ifcase#1\or\faStar\faStar[regular]\faStar[regular]\or\faStar\faStar\faStar[regular]\or\faStar\faStar\faStar\fi}}}

% ---------- capçalera i peu de les pàgines ----------
% primera pàgina: logos institucionals al peu
\fancypagestyle{primera}{\fancyhf{}\renewcommand{\headrulewidth}{0pt}
  \fancyfoot[L]{\raisebox{-0.32\height}{\includegraphics[height=0.78cm]{logo-gva.png}}\hspace{5mm}%
    \raisebox{-0.32\height}{\includegraphics[height=0.78cm]{logo-vermellar.png}}}
  \fancyfoot[R]{\scriptsize\color{muted}\href{https://creativecommons.org/licenses/by-nc-sa/4.0/deed.ca}{\ccbyncsa}\ \ CC BY-NC-SA 4.0 · fentmates.org\qquad\small\thepage}}
% la resta: el símbol a l'esquerra i el curs a la dreta, amb el títol al centre com abans
\fancyhead[L]{\raisebox{-0.3\height}{\includegraphics[height=0.62cm]{\fm@marca}}}
\fancyhead[R]{\small\bfseries\color{cursfosc}\fm@cursnom}
\renewcommand{\headrule}{\color{curs!45}\hrule height \headrulewidth}
\makeatother

%   \fmcurs{1eso}
% i, al començament del document:
%   \capcalera{Matemàtiques}{1r ESO}{Fitxa d'exercicis · Unitat 2}{Divisibilitat}
% Opcional abans de \input: \def\fmneed{14} (espai mínim abans de cada secció).
% =====================================================================
\makeatletter
\def\fm@logodir{../../logo/pdf}
\providecommand{\fmneed}{10}
\def\fm@cursnom{}

% clau del curs, color i color del text damunt del color
\newcommand{\fm@curs}[3]{\@namedef{fm@color@#1}{#2}\@namedef{fm@text@#1}{#3}}
\fm@curs{1eso}{E6B84A}{navy}
\fm@curs{2eso}{3B82F6}{white}
\fm@curs{3eso}{FF7A50}{white}
\fm@curs{4eso-a}{A78BFA}{white}
\fm@curs{4eso-b}{6D28D9}{white}
\fm@curs{1batx-mi}{9B1C1C}{white}
\fm@curs{1batx-ccss}{EF6B6B}{white}
\fm@curs{2batx-mii}{1E6B47}{white}
\fm@curs{2batx-ccss}{2BB3A3}{white}

\newcommand{\fmcurs}[1]{%
  \@ifundefined{fm@color@#1}{\PackageError{fentmates}{Curs desconegut: #1}{}}{}%
  \edef\fm@tmp{\noexpand\definecolor{curs}{HTML}{\@nameuse{fm@color@#1}}%
               \noexpand\colorlet{curstext}{\@nameuse{fm@text@#1}}}\fm@tmp
  \colorlet{cursfosc}{curs!72!black}%
  \def\fm@logo{\fm@logodir/vertical-#1.pdf}%
  \def\fm@marca{\fm@logodir/marca-#1.pdf}}

% ---------- capçalera de la primera pàgina ----------
% #1 matèria (petit), #2 curs, #3 tipus de document i unitat, #4 títol
\newcommand{\capcalera}[4]{%
  \gdef\fm@cursnom{#2}%
  \par\noindent
  \begin{tcolorbox}[colback=curs,colframe=curs,boxrule=0pt,arc=2.2mm,width=3.5cm,height=2.15cm,
      valign=center,left=3mm,right=2mm,top=1mm,bottom=1mm,nobeforeafter,box align=center]
    \color{curstext}{\footnotesize #1}\par\vspace{1pt}{\fontsize{19}{21}\selectfont\bfseries #2\par}%
  \end{tcolorbox}\hspace{5mm}%
  \begin{minipage}[c]{\dimexpr\linewidth-3.5cm-5mm-2.9cm\relax}
    {\color{cursfosc}\small\bfseries\MakeUppercase{#3}}\par\vspace{2pt}
    {\color{navy}\fontsize{20}{23}\selectfont\bfseries #4\par}
  \end{minipage}\hfill
  \begin{minipage}[c]{2.6cm}\raggedleft\includegraphics[width=2.6cm]{\fm@logo}\end{minipage}\par
  \vspace{2pt}}

% ---------- barres de secció del color del curs ----------
\renewcommand{\seccio}[2]{\par\needspace{\fmneed\baselineskip}\bigskip
  \begin{tcolorbox}[colback=curs,colframe=curs,coltext=curstext,arc=2mm,boxrule=0pt,left=3mm,right=3mm,top=1.2mm,bottom=1.2mm]
  {\large\bfseries #1\quad #2}
  \end{tcolorbox}\vspace{-2pt}}

% ---------- estrelles de dificultat del color del curs (a Batxillerat no se'n posen) ----------
\@ifundefined{estrelles}{}{%
  \renewcommand{\estrelles}[1]{{\color{curs}\ifcase#1\or\faStar\faStar[regular]\faStar[regular]\or\faStar\faStar\faStar[regular]\or\faStar\faStar\faStar\fi}}}

% ---------- capçalera i peu de les pàgines ----------
% primera pàgina: logos institucionals al peu
\fancypagestyle{primera}{\fancyhf{}\renewcommand{\headrulewidth}{0pt}
  \fancyfoot[L]{\raisebox{-0.32\height}{\includegraphics[height=0.78cm]{logo-gva.png}}\hspace{5mm}%
    \raisebox{-0.32\height}{\includegraphics[height=0.78cm]{logo-vermellar.png}}}
  \fancyfoot[R]{\scriptsize\color{muted}\href{https://creativecommons.org/licenses/by-nc-sa/4.0/deed.ca}{\ccbyncsa}\ \ CC BY-NC-SA 4.0 · fentmates.org\qquad\small\thepage}}
% la resta: el símbol a l'esquerra i el curs a la dreta, amb el títol al centre com abans
\fancyhead[L]{\raisebox{-0.3\height}{\includegraphics[height=0.62cm]{\fm@marca}}}
\fancyhead[R]{\small\bfseries\color{cursfosc}\fm@cursnom}
\renewcommand{\headrule}{\color{curs!45}\hrule height \headrulewidth}
\makeatother

\fmcurs{2batx-mii}

\begin{document}
\thispagestyle{primera}

\capcalera{Matemáticas II}{2\textsuperscript{o} Bach.}{Ficha de ejercicios · Unidad 3}{\raggedright Vectores y aplicaciones. Rectas, planos y haces de planos}
\vspace{10pt}
Nombre y apellidos: \rule{6.6cm}{0.4pt}\hfill Grupo: \rule{1.6cm}{0.4pt}\hfill Fecha: \rule{2.2cm}{0.4pt}
\vspace{8pt}

\begin{tcolorbox}[colback=soft,colframe=soft,arc=2mm,left=3mm,right=3mm,top=1.5mm,bottom=1.5mm]
\small Resuelve los ejercicios en el cuaderno justificando cada paso. Haz un esbozo cuando te ayude a ver la situación y comprueba siempre los resultados (que el punto hallado está en los planos, que el vector normal es perpendicular a los vectores directores\ldots). Los ejercicios marcados con {\color{teal}\bfseries\faUniversity\ PAU} son de las pruebas de acceso a la universidad de otras comunidades autónomas (2024--2026); se han quitado los apartados de ángulos y distancias, que son de la unidad siguiente.
\end{tcolorbox}

% =====================================================
\seccio{1}{Vectores en el espacio y aplicaciones}

\ex{1} Sean $A(1,\,-2,\,3)$, $B(3,\,-1,\,1)$ y $C(-1,\,2,\,7)$.\begin{tasks}(2)
  \task Calcula $\overrightarrow{AB}$, $\overrightarrow{AC}$ y sus módulos.
  \task Halla un vector unitario con la dirección y el sentido de $\overrightarrow{AC}$.
  \task Halla un vector de módulo 9 paralelo a $\overrightarrow{AB}$ y de sentido contrario.
  \task Calcula $3\overrightarrow{AB}+\overrightarrow{AC}$.
\end{tasks}

\ex{1} El punto medio del segmento $AB$ es $M(1,\,-2,\,4)$ y $A(3,\,0,\,-1)$. Calcula $B$. Halla también el simétrico de $P(2,\,-3,\,1)$ respecto del punto $Q(0,\,1,\,-2)$.\recorda{El punto medio de $A$ y $B$ es $M=\frac{A+B}{2}$; por tanto, $B=2M-A$. El simétrico de un punto respecto de otro se calcula igual.}

\ex{2} Divide en tres partes iguales el segmento de extremos $A(-2,\,1,\,5)$ y $B(4,\,-5,\,-1)$.

\ex{2} Halla $a$ y $b$ para que los puntos $A(1,\,0,\,2)$, $B(2,\,1,\,0)$ y $C(a,\,3,\,b)$ estén alineados.

\ex{2} Los puntos $A(1,\,2,\,-1)$, $B(3,\,0,\,2)$ y $C(2,\,-1,\,4)$ son tres vértices consecutivos del paralelogramo $ABCD$. Calcula el cuarto vértice $D$ y el centro del paralelogramo.

\ex{2} El baricentro del triángulo $ABC$ es $G(1,\,2,\,0)$, con $A(2,\,1,\,-3)$ y $B(-1,\,4,\,2)$. Halla el vértice $C$.

\ex{2} Halla el punto $P$ del segmento $AB$, con $A(1,\,-1,\,2)$ y $B(7,\,5,\,-4)$, que cumple $\overline{AP}=2\,\overline{PB}$.

\ex{2} \pau{Castilla-La Mancha, extraordinaria 2026} Se consideran los puntos $A(1,\,2,\,-1)$, $B(3,\,6,\,1)$ y $C(m,\,m,\,m-2)$, donde $m$ es un número real. Determina el valor de $m$ para que los puntos $A$, $B$ y $C$ estén alineados.

\ex{2} \pau{Cataluña, extraordinaria 2026} Ana quiere hacer un croquis en 3D de una caja con un programa informático. Sitúa un lado de la tapa desde el vértice $A(1,\,1,\,1)$ hasta el vértice $B(4,\,1,\,5)$, y el otro lado de la tapa desde el vértice $A$ hasta el vértice $C(1,\,10,\,1)$. ¿Cuáles son las coordenadas del cuarto vértice de la tapa? Calcula las longitudes de los dos lados de la tapa.


% =====================================================
\seccio{2}{Dependencia lineal y bases}

\ex{1} Estudia si los vectores $\vec u=(1,\,2,\,-1)$, $\vec v=(2,\,1,\,3)$ y $\vec w=(0,\,3,\,-5)$ son linealmente independientes. Si no lo son, escribe $\vec w$ como combinación lineal de $\vec u$ y $\vec v$.

\ex{2} Considera los vectores $\vec u=(1,\,1,\,1)$, $\vec v=(1,\,a,\,1)$ y $\vec w=(1,1,a^2)$.\begin{tasks}(1)
  \task ¿Para qué valores de $a$ forman una base de $\mathbb R^3$?
  \task Calcula el rango del conjunto $\{\vec u,\vec v,\vec w\}$ para los demás valores de $a$.
\end{tasks}

\ex{2} Comprueba que $\mathcal B=\{(1,\,1,\,0),(0,\,1,\,1),(1,\,0,\,1)\}$ es una base y calcula las coordenadas del vector $\vec x=(3,\,1,\,2)$ en esta base.

\ex{2} ¿Se puede escribir $(1,\,2,\,3)$ como combinación lineal de $(1,\,1,\,0)$ y $(0,\,1,\,1)$? ¿Y $(2,\,5,\,3)$?

\ex{3} Sea $\{\vec u,\vec v,\vec w\}$ una base. Estudia si son linealmente independientes:\begin{tasks}(2)
  \task $\{\vec u+\vec v,\ \vec v+\vec w,\ \vec u+\vec w\}$
  \task $\{\vec u-\vec v,\ \vec v-\vec w,\ \vec w-\vec u\}$
\end{tasks}\recorda{Escribe cada vector nuevo con coordenadas en la base $\{\vec u,\vec v,\vec w\}$ y calcula el determinante de estas coordenadas.}


% =====================================================
\seccio{3}{Producto escalar}

\ex{1} Dados $\vec u=(1,\,-2,\,2)$ y $\vec v=(3,\,0,\,-4)$, calcula:\begin{tasks}(2)
  \task $\vec u\cdot\vec v$, $|\vec u|$ y $|\vec v|$.
  \task El ángulo que forman.
  \task La proyección escalar de $\vec u$ sobre $\vec v$.
  \task El vector proyección de $\vec u$ sobre $\vec v$.
\end{tasks}

\ex{2} Halla los vectores de módulo $\sqrt{29}$ que son perpendiculares a la vez a $(1,\,2,\,-1)$ y a $(2,\,0,\,1)$ (sin utilizar el producto vectorial).

\ex{2} Halla $m$ para que los vectores $(2,\,m,\,1)$ y $(1,\,1,\,2)$ formen un ángulo de $60^\circ$.

\ex{2} Calcula los ángulos que forma el vector $\vec v=(2,\,-1,\,2)$ con los ejes de coordenadas. ¿Puede haber un vector que forme $45^\circ$ con el eje $OX$ y $60^\circ$ con el eje $OY$? ¿Qué ángulo formaría con el eje $OZ$?

\ex{3} Sabemos que $|\vec u|=3$, $|\vec v|=4$ y $\vec u\cdot\vec v=6$. Calcula:\begin{tasks}(2)
  \task El ángulo que forman $\vec u$ y $\vec v$.
  \task $|\vec u+\vec v|$ y $|\vec u-\vec v|$.
  \task $(2\vec u-\vec v)\cdot(\vec u+3\vec v)$. ¿Qué observas?
\end{tasks}\recorda{$|\vec a|^2=\vec a\cdot\vec a$, y el producto escalar es distributivo.}

\ex{2} Demuestra que el triángulo de vértices $A(1,\,1,\,1)$, $B(3,\,2,\,-1)$ y $C(2,\,3,\,3)$ es rectángulo e isósceles, y calcula sus ángulos.

\ex{1} Halla $m$ para que los vectores $(m,\,1,\,3)$ y $(2,\,m,\,-1)$ sean perpendiculares.


% =====================================================
\seccio{4}{Producto vectorial y producto mixto. Áreas y volúmenes}

\ex{1} Sean $\vec u=(2,\,-1,\,3)$ y $\vec v=(1,\,0,\,-2)$. Calcula $\vec u\times\vec v$ y $\vec v\times\vec u$, y comprueba que $\vec u\times\vec v$ es perpendicular a $\vec u$ y a $\vec v$.

\ex{2} Dados $A(1,\,0,\,1)$, $B(2,\,1,\,3)$ y $C(0,\,2,\,1)$, halla los vectores unitarios perpendiculares al plano que determinan y el área del triángulo $ABC$.

\ex{2} Los puntos $A(1,\,1,\,0)$, $B(3,\,2,\,1)$ y $D(0,\,3,\,2)$ son vértices del paralelogramo $ABCD$. Calcula el vértice $C$ y el área del paralelogramo.

\ex{2} Halla $k$ para que el triángulo de vértices $A(1,\,0,\,0)$, $B(0,\,1,\,0)$ y $C(0,\,0,\,k)$ tenga área $\tfrac32$.

\ex{2} Halla $a$ para que el tetraedro de vértices $A(1,\,1,\,1)$, $B(3,\,1,\,0)$, $C(1,\,4,\,2)$ y $D(1,\,a,\,3)$ tenga volumen 2.

\ex{2} Halla $a$ para que los puntos $A(1,\,2,\,1)$, $B(2,\,3,\,0)$, $C(1,\,0,\,3)$ y $D(a,\,1,\,a)$ sean coplanarios.

\ex{3} Sabemos que $[\vec u,\vec v,\vec w]=5$. Calcula:\begin{tasks}(3)
  \task $[\vec w,\vec v,\vec u]$
  \task $[2\vec u,\ \vec v-\vec w,\ 3\vec w]$
  \task $[\vec u+\vec v,\ \vec v,\ \vec w]$
\end{tasks}\recorda{El producto mixto es un determinante: aplica sus propiedades (permutar filas, sacar factores, sumar a una fila una combinación de las otras).}

\ex{2} \pau{Andalucía, ordinaria 2026} Considera los puntos $A(1,\,2,\,0)$, $B(2,\,m,\,1)$, $C(3,\,4,\,2)$ y $D(1,\,-1,\,m)$.\begin{tasks}(1)
  \task Halla los valores de $m$ para los que los puntos son coplanarios.
  \task Para $m=1$, calcula el área del triángulo de vértices $A$, $B$, $C$ y el volumen del tetraedro de vértices $A$, $B$, $C$, $D$.
\end{tasks}

\ex{2} \pau{Aragón, ordinaria 2024} Dados los puntos $P_1(-2,\,1,\,1)$, $P_2(0,\,a,\,-2)$, $P_3(-1,\,1,\,-1)$ y $P_4(1,\,3,\,-3)$:\begin{tasks}(1)
  \task Calcula los valores de $a$ para que el tetraedro de vértices $P_1$, $P_2$, $P_3$ y $P_4$ tenga volumen $\tfrac13$.
  \task Calcula el valor de $a$ para que los cuatro puntos sean coplanarios.
\end{tasks}

\ex{2} \pau{Asturias, ordinaria 2025} Se consideran los puntos $A(1,\,2,\,3)$, $B(-2,\,1,\,4)$, $C(3,\,0,\,5)$ y $D(0,\,-1,\,2)$.\begin{tasks}(1)
  \task Estudia si los cuatro puntos pertenecen a un mismo plano.
  \task Calcula el área del triángulo de vértices $A$, $B$ y $C$.
  \task Calcula el volumen del tetraedro que forman los cuatro puntos.
\end{tasks}

\ex{2} \pau{Castilla-La Mancha, ordinaria 2026} Dados los vectores $\vec u=(1,\,-1,\,a)$, $\vec v=(2,\,a,\,1)$ y $\vec w=(-1,\,1,\,-1)$:\begin{tasks}(1)
  \task Halla los valores reales de $a$ para que los tres vectores sean coplanarios. Justifica la respuesta.
  \task Halla los valores reales de $a$ para que el volumen del paralelepípedo formado por los tres vectores sea $10\ u^3$.
\end{tasks}

\ex{2} \pau{Madrid, extraordinaria 2025} Sean los puntos $A(1,\,1,\,2)$, $B(2,\,-1,\,0)$, $C(-2,\,0,\,3)$ y $D(2,\,-3,\,-1)$.\begin{tasks}(1)
  \task Comprueba que los puntos no son coplanarios y calcula el volumen del tetraedro que determinan.
  \task Calcula el área de la cara $ABC$ del tetraedro y la longitud de la altura del tetraedro que parte del vértice $D$.
\end{tasks}\recorda{El volumen de una pirámide es un tercio del área de la base por la altura.}

\ex{2} \pau{Navarra, ordinaria 2024} Los puntos $A(4,\,-2,\,-3)$, $B(2,\,-1,\,1)$ y $C(0,\,-3,\,-1)$ son vértices de un rombo.\begin{tasks}(1)
  \task Halla el cuarto vértice del rombo.
  \task Calcula el área del rombo.
\end{tasks}

\ex{2} \pau{Aragón, extraordinaria 2024} Sean $A(1,\,2,\,3)$, $B(1,\,0,\,-1)$ y $C(2,\,2,\,2)$ tres puntos del espacio, $\vec v_1$ el vector que va de $A$ a $B$, $\vec v_2$ el vector que va de $B$ a $C$ y $\vec v_3$ el vector que va de $C$ a $A$.\begin{tasks}(1)
  \task Estudia si los vectores $\vec v_1$, $\vec v_2$ y $\vec v_3$ son linealmente independientes.
  \task Calcula el área del triángulo de vértices $A$, $B$ y $C$.
\end{tasks}


% =====================================================
\seccio{5}{Ecuaciones de rectas y planos}

\ex{1} Escribe las ecuaciones vectorial, paramétricas, continua e implícitas de la recta $r$ que pasa por $A(2,\,-1,\,3)$ y $B(4,\,0,\,1)$. ¿Los puntos $C(0,\,-2,\,5)$ y $D(6,\,1,\,0)$ son de $r$?

\ex{1} Escribe en forma paramétrica la recta $r:\left\{\begin{aligned}x-y+z&=2\\2x+y-z&=1\end{aligned}\right.$

\ex{1} Halla las ecuaciones paramétricas y general del plano que pasa por $A(1,\,0,\,2)$, $B(2,\,1,\,0)$ y $C(0,\,3,\,1)$.

\ex{2} Halla la ecuación general del plano que pasa por $P(2,\,-1,\,1)$ y es paralelo a las rectas \mbox{$r:\dfrac{x-1}{2}=\dfrac y1=\dfrac{z+1}{-1}$} y $s:(x,\,y,\,z)=(0,\,1,\,0)+\mu(1,\,0,\,3)$.

\ex{2} Comprueba que las rectas \mbox{$r:\dfrac{x-1}{1}=\dfrac{y+1}{2}=\dfrac{z}{-1}$} y \mbox{$s:\dfrac{x+1}{-2}=\dfrac{y-2}{-4}=\dfrac{z-1}{2}$} son paralelas y halla el plano que las contiene.

\ex{2} Escribe la ecuación segmentaria y la general del plano que corta los ejes en $(2,\,0,\,0)$, $(0,\,-3,\,0)$ y $(0,\,0,\,6)$. ¿Qué volumen tiene el tetraedro que forman el plano $4x-2y+3z=12$ y los planos coordenados?

\ex{2} Halla el plano que pasa por el origen y es perpendicular a la recta $r:\left\{\begin{aligned}x+y&=1\\y-z&=2\end{aligned}\right.$, y la recta que pasa por $A(1,\,2,\,3)$ y es perpendicular al plano $2x-y+3z=5$.

\ex{1} Pasa a forma general el plano $(x,\,y,\,z)=(1,\,0,\,-1)+\lambda(1,\,2,\,0)+\mu(0,\,1,\,3)$ y halla $a$, $b$ para que $P(a,\,1,\,b)$ sea de la recta \mbox{$\dfrac{x-2}{3}=\dfrac{y+1}{2}=\dfrac{z-1}{-1}$}.

\ex{1} \pau{Andalucía, ordinaria 2025} Considera el punto $P(1,\,1,\,1)$ y la recta \mbox{$r:\dfrac{x-1}{1}=\dfrac{y-2}{2}=\dfrac{z-3}{2}$}. Halla el plano que pasa por el punto $P$ y contiene a la recta $r$.

\ex{2} \pau{Aragón, ordinaria 2025}\begin{tasks}(1)
  \task Halla las ecuaciones paramétricas de la recta $s$ que pasa por el punto $P(4,\,-3,\,0)$ y es perpendicular al plano $\pi:x-2y+z-1=0$.
  \task Halla la ecuación del plano que contiene el punto $Q(1,\,2,\,3)$ y la recta $r:\left\{\begin{aligned}x+y+z&=1\\x-2y&=0\end{aligned}\right.$
\end{tasks}

\ex{2} \pau{Cantabria, extraordinaria 2024} Sean $A(0,\,3,\,2)$, $B(4,\,1,\,3)$, $C(2,\,3,\,4)$ y $D(0,\,1,\,2)$ los vértices de un tetraedro.\begin{tasks}(1)
  \task Obtén la ecuación vectorial del plano determinado por los puntos $A$, $B$ y $C$.
  \task Calcula el volumen del tetraedro.
\end{tasks}

\ex{2} \pau{Castilla-La Mancha, extraordinaria 2026} Se consideran los puntos $A(1,\,1,\,0)$, $B(2,\,0,\,1)$, $C(0,\,2,\,1)$ y $D(a,\,1,\,2)$, con $a$ un número real.\begin{tasks}(1)
  \task Determina la ecuación del plano que pasa por los puntos $A$, $B$ y $C$.
  \task Determina los valores de $a$ para que el tetraedro formado por los cuatro puntos tenga volumen 3 unidades cúbicas.
\end{tasks}

\ex{2} \pau{Castilla y León, extraordinaria 2026} Los puntos $A(1,\,1,\,1)$ y $B(0,\,2,\,0)$ son vértices consecutivos del paralelogramo $ABCD$, que tiene el centro en $E(0,\,0,\,1)$.\begin{tasks}(1)
  \task Calcula las coordenadas de los otros vértices.
  \task Halla la ecuación general del plano que contiene el paralelogramo.
  \task Halla el volumen del tetraedro de vértices $A$, $B$, $E$ y el origen de coordenadas $O$.
\end{tasks}

\ex{2} \pau{Cataluña, ordinaria 2026} Considera los puntos $P(1,\,0,\,-1)$, $Q(3,\,-2,\,0)$ y $R(1,\,1,\,1)$.\begin{tasks}(1)
  \task Calcula la ecuación del plano que contiene los puntos $P$, $Q$ y $R$.
  \task Comprueba que el área del triángulo $PQR$ es $\tfrac{3\sqrt5}{2}\ u^2$.
  \task Determina las condiciones que deben cumplir las coordenadas de un cuarto punto $S(x,\,y,\,z)$ para que $P$, $Q$, $R$ y $S$ formen un tetraedro de volumen 1.
\end{tasks}

\ex{2} \pau{Extremadura, extraordinaria 2026} Dados los puntos $A(4,\,3,\,3)$, $B(3,\,3,\,4)$, $C(1,\,0,\,4)$ y $D(4,\,0,\,1)$:\begin{tasks}(1)
  \task Justifica de forma razonada que los puntos $A$, $B$, $C$ y $D$ están en el mismo plano y halla la ecuación del plano que determinan.
  \task Calcula el valor de $k$ para que el punto $P(7,\,k,\,2)$ pertenezca a la recta $r$ que pasa por los puntos $A$ y $C$.
\end{tasks}

\ex{2} \pau{Extremadura, ordinaria 2026}\begin{tasks}(1)
  \task Dados los vectores $\vec u=(0,\,3,\,-4)$ y $\vec v=(12,\,4,\,3)$, comprueba que el paralelogramo que forman es un rectángulo y calcula su área.
  \task Halla la ecuación del plano formado por los vectores $\vec u$ y $\vec v$ y el punto $(1,\,0,\,0)$.
  \task Obtén un vector perpendicular a los vectores $\vec u$ y $\vec v$ de módulo 13.
\end{tasks}

\ex{2} \pau{Madrid, ordinaria 2025} Dados la recta \mbox{$r:\dfrac{x-1}{2}=\dfrac y0=\dfrac{z-2}{1}$} y el plano $\pi:x+2y-3z=1$, halla una ecuación del plano que contiene a $r$ y es perpendicular a $\pi$.


% =====================================================
\seccio{6}{Posiciones relativas de rectas y planos}

\ex{1} Estudia la posición relativa de \mbox{$r:\dfrac{x-1}{2}=\dfrac{y+2}{-1}=\dfrac z3$} y $s:(x,\,y,\,z)=(3,\,-3,\,4)+\mu(4,\,-2,\,6)$. Si son coplanarias, halla el plano que las contiene.

\ex{1} Estudia la posición relativa de $r:(x,\,y,\,z)=(1,\,0,\,2)+\lambda(2,\,1,\,-1)$ y $s:(x,\,y,\,z)=(0,\,3,\,1)+\mu(1,\,-2,\,1)$.

\ex{2} Halla $a$ para que las rectas \mbox{$r:\dfrac{x-1}{1}=\dfrac{y-2}{2}=\dfrac{z-a}{1}$} y $s:(x,\,y,\,z)=(0,\,1,\,3)+\mu(1,\,0,\,1)$ se corten, y calcula el punto de corte.

\ex{2} Estudia, según los valores de $k$, la posición relativa de $r:(x,\,y,\,z)=(2,\,-1,\,0)+\lambda(1,\,1,\,2)$ y $\pi:3x-y-z+k=0$.

\ex{2} Estudia, según los valores de $m$, la posición de la recta $r:\left\{\begin{aligned}x+y-z&=1\\2x-y+z&=2\end{aligned}\right.$ y el plano $\pi:x+my+z=2$. Halla el punto de corte cuando $m=1$.

\ex{2} Estudia la posición relativa de los planos $\pi_1:2x-y+3z=1$ y $\pi_2:4x+ay+6z=b$ según los valores de $a$ y $b$. Para $a=1$, $b=2$, halla la recta en la que se cortan.

\ex{3} Halla la recta que pasa por $P(1,\,1,\,1)$, es paralela al plano $\pi:x+y+z=0$ y corta a la recta $s:(x,\,y,\,z)=(2,\,1,\,0)+\mu(0,\,1,\,1)$.\recorda{La recta buscada está contenida en el plano paralelo a $\pi$ que pasa por $P$. Corta a $s$ en el punto donde $s$ corta a este plano.}

\ex{2} \pau{Asturias, extraordinaria 2024} Se consideran los puntos $A(1,\,1,\,1)$, $B(1,\,0,\,2)$, $C(-1,\,1,\,3)$ y $D(-1,\,0,\,1)$.\begin{tasks}(1)
  \task Estudia si hay un plano que contenga los cuatro puntos.
  \task Calcula la recta $r$ que pasa por $D$ y es perpendicular al plano $\pi$ que contiene $A$, $B$ y $C$.
  \task Calcula el punto $P$ de intersección de $r$ y $\pi$.
\end{tasks}

\ex{2} \pau{Canarias, extraordinaria 2025} Se tienen las rectas $r:\left\{\begin{aligned}x+y-2z&=3\\2x+y-6z&=2\end{aligned}\right.$ \ y \ \mbox{$s:x-1=\dfrac y3=\dfrac{z+1}{2}$}.\begin{tasks}(1)
  \task Comprueba que $r$ y $s$ son coplanarias.
  \task Halla la ecuación del plano que las contiene.
\end{tasks}

\ex{2} \pau{Canarias, ordinaria 2026} Se conocen el plano $\pi:2x+3y=0$ y el punto $P(1,\,-5,\,1)$.\begin{tasks}(1)
  \task Halla la recta $r$ perpendicular al plano $\pi$ que pasa por el punto $P$.
  \task Halla la posición relativa de la recta $s:\left\{\begin{aligned}-x+3y-3z&=-3\\4x+9y-2z&=-2\end{aligned}\right.$ \ con el plano $\pi$.
\end{tasks}

\ex{3} \pau{Canarias, extraordinaria 2024} Considera las rectas $r_1:\left\{\begin{aligned}x-3y+2z+2&=0\\2x+y-3z&=3\end{aligned}\right.$ \ y \ \mbox{$r_2:\dfrac{1-x}{2}=y=\dfrac{1-z}{2}$}.\begin{tasks}(1)
  \task Estudia la posición relativa de las dos rectas.
  \task Halla la ecuación de la recta $s$ que tiene dirección perpendicular a las dos rectas y pasa por $P\left(0,\,\tfrac12,\,0\right)$. Calcula el punto de corte de $s$ con $r_1$.
\end{tasks}

\ex{2} \pau{Cantabria, ordinaria 2024} Considera la recta $r:\left\{\begin{aligned}x+y+z+5&=0\\x+2y-z&=0\end{aligned}\right.$ \ y el plano $\pi:2x+y-az=3$. Razona si es posible dar algún valor al parámetro $a$ para que:\begin{tasks}(1)
  \task la recta esté contenida en el plano;
  \task la recta y el plano sean paralelos;
  \task la recta y el plano se corten. En caso afirmativo, da un valor de $a$ y calcula dónde se cortan.
\end{tasks}

\ex{2} \pau{Castilla y León, ordinaria 2025} Sean las rectas \mbox{$r:\dfrac x4=\dfrac{y-1}{-2}=\dfrac z2$} y \mbox{$s:x-1=\dfrac{y-m}{m-1}=\dfrac{z-3}{3}$}.\begin{tasks}(1)
  \task Comprueba que las rectas $r$ y $s$ se cortan para cualquier valor de $m$.
  \task Para $m=6$, halla el punto de intersección de las rectas $r$ y $s$.
\end{tasks}

\ex{2} \pau{Castilla y León, extraordinaria 2025} Sean el plano $\pi:5x+my+z=2$ y la recta $r:(x,\,y,\,z)=(1,\,1,\,0)+t(-1,\,-1,\,2)$.\begin{tasks}(1)
  \task Determina la posición relativa de $r$ y $\pi$ en función de $m$.
  \task Para $m=1$, calcula el plano $\pi'$ que contiene a $r$ y es perpendicular a $\pi$.
\end{tasks}

\ex{1} \pau{Galicia, extraordinaria 2025} Dadas las rectas\[r:\dfrac{x-1}{2}=\dfrac{y-2}{-1}=\dfrac{z-1}{1}\qquad\text{y}\qquad s:\dfrac{x-2}{4}=\dfrac{y-1}{-2}=\dfrac{z-1}{2}\]\begin{tasks}(1)
  \task Calcula la posición relativa de las rectas $r$ y $s$.
  \task Obtén la ecuación del plano que contiene las rectas $r$ y $s$.
\end{tasks}

\ex{2} \pau{Galicia, extraordinaria 2024} Considera el plano $\pi:x+2y-2z=0$ y la recta $r$ que pasa por los puntos $A(2,\,1,\,2)$ y $B(0,\,1,\,1)$.\begin{tasks}(1)
  \task Estudia la posición relativa de la recta $r$ y el plano $\pi$.
  \task Obtén la ecuación implícita o general del plano que contiene a $r$ y es perpendicular a $\pi$.
\end{tasks}

\ex{1} \pau{Galicia, ordinaria 2024} Considera el plano $\pi:4x+2y+bz=2$ y la recta \mbox{$r:\dfrac{x-2}{3}=\dfrac{y-c}{2}=\dfrac{z-3}{4}$}, donde $b$ y $c$ son parámetros reales. Calcula los valores de $b$ y $c$ para que la recta $r$ esté contenida en $\pi$.

\ex{2} \pau{Murcia, extraordinaria 2025} Considera la recta $r:\left\{\begin{aligned}x+y+z&=1\\-x-2y+z&=0\end{aligned}\right.$ \ y el plano $\pi:2x+y+az=b$.\begin{tasks}(1)
  \task Determina los valores de $a$ y $b$ para que el plano contenga a la recta.
  \task ¿Para qué valores de $a$ y $b$ la recta corta al plano? Halla el punto de corte en el caso $a=0$, $b=7$.
\end{tasks}

\ex{2} \pau{Navarra, extraordinaria 2024} Se consideran el plano $\pi:2x+y-z-5=0$, la recta $r:\left\{\begin{aligned}x+2z+3&=0\\-x-y+z+4&=0\end{aligned}\right.$ \ y los puntos $A(3,\,2,\,-1)$ y $B(1,\,1,\,-1)$. Sea $C$ la intersección de la recta y el plano.\begin{tasks}(1)
  \task Demuestra que los puntos $A$, $B$ y $C$ no están alineados.
  \task Calcula el área del triángulo que forman los tres puntos.
\end{tasks}

\ex{3} \pau{Navarra, ordinaria 2025} Calcula la ecuación continua de la recta $t$ que pasa por el punto $P(2,\,0,\,-1)$ y corta a las rectas $s:\left\{\begin{aligned}2x+y-3z-6&=0\\2x-3z-8&=0\end{aligned}\right.$ \ y \ \mbox{$r:\dfrac{x+1}{2}=\dfrac y1=\dfrac{z+2}{1}$}.\recorda{La recta $t$ es la intersección del plano que contiene $P$ y $s$ con el plano que contiene $P$ y $r$.}

\ex{2} \pau{País Vasco, extraordinaria 2024} Se consideran la recta $r:\left\{\begin{aligned}2x-y+z&=0\\x-y+4z&=1\end{aligned}\right.$ \ y el plano $\pi:2x-3y+Az=10$.\begin{tasks}(1)
  \task Calcula el valor del parámetro $A$ para que la recta $r$ y el plano $\pi$ sean paralelos.
  \task Si $A=21$, calcula la intersección del plano $\pi$ y la recta $r$.
\end{tasks}

\ex{1} \pau{Aragón, ordinaria 2026} Sean los planos $\pi:x-2y+4z+1=0$ y $\pi_d:2x-3y+6z=d$, con $d\in\mathbb R$.\begin{tasks}(1)
  \task ¿Hay algún valor de $d$ para el que los planos $\pi$ y $\pi_d$ no se cortan?
  \task Calcula, para $d=-2$, los puntos de corte entre los planos $\pi$ y $\pi_d$.
\end{tasks}


% =====================================================
\seccio{7}{Posición relativa de tres planos}

\ex{2} Estudia, según el valor de $a$, la posición relativa de los planos $\pi_1:x+y-z=1$, $\pi_2:2x-y+z=2$ y $\pi_3:x-2y+2z=a$. Cuando tengan una recta común, hállala.

\ex{3} Discute, según los valores de $k$, la posición relativa de los planos $kx+y+z=k$, \ $x+ky+z=k$, \ $x+y+kz=k$, y halla el punto común cuando haya uno solo.

\ex{2} Halla $a$ y $b$ para que los planos $x+y+z=2$, $x-y+2z=1$ y $3x+y+az=b$ tengan una recta en común. ¿Qué pasa si $a=4$ y $b\neq5$? ¿Y si $a\neq4$?

\ex{2} Dados $\pi_1:x-y=0$ y $\pi_2:y-z=0$, escribe un plano $\pi_3$ tal que los tres planos\begin{tasks}(2)
  \task formen un prisma;
  \task se corten en un solo punto.
\end{tasks}

\ex{3} \pau{La Rioja, extraordinaria 2026} Determina los valores de $a$ para que los planos $\pi_1:x+y+z=1$, \ $\pi_2:(a-1)x+y+z=a$ \ y \ $\pi_3:x+(a-1)y-z=0$\begin{tasks}(1)
  \task se corten en un punto. En este caso, determina el punto de corte.
  \task se corten en una recta. En este caso, determina la recta en forma paramétrica.
\end{tasks}


% =====================================================
\seccio{8}{Haces de planos}

\ex{1} Sea $r:\left\{\begin{aligned}x-y+z&=1\\2x+y&=3\end{aligned}\right.$. Escribe el haz de planos de arista $r$ y halla el plano del haz que pasa por el origen.

\ex{2} Con la recta $r$ del ejercicio anterior, halla el plano que contiene a $r$ y es paralelo al eje $OZ$, y el plano que contiene a $r$ y es perpendicular a $\pi:x+2y-z=0$.

\ex{2} Halla el plano paralelo a $2x-y+2z=1$ que pasa por $A(1,\,-1,\,2)$, y los planos de este haz que forman con los planos coordenados un tetraedro de volumen 9.

\ex{2} Escribe la recta \mbox{$r:\dfrac{x-1}{2}=\dfrac y1=\dfrac{z+1}{3}$} como intersección de dos planos y halla, con el haz de planos, el plano que contiene a $r$ y pasa por $P(1,\,1,\,1)$.

\ex{2} \pau{Murcia, extraordinaria 2026} Determina la ecuación del plano que contiene el punto $P(-2,\,2,\,-2)$ y la recta $r$ dada como intersección de dos planos: \ $r:\left\{\begin{aligned}2x+y-z&=4\\x+y-2z&=3\end{aligned}\right.$

\ex{2} \pau{País Vasco, ordinaria 2025} Se consideran las rectas $r_1:\left\{\begin{aligned}x+y-2z&=0\\2x-3y+z&=1\end{aligned}\right.$ \ y \ $r_2:\left\{\begin{aligned}x&=3t\\y&=1-2t\\z&=2+t\end{aligned}\right.$ \ Calcula la ecuación del plano que contiene a $r_1$ y pasa por el punto de intersección del plano $\pi:x-3y-2z+7=0$ y la recta $r_2$.


% =====================================================
\seccio{9}{Problemas y cuestiones}

\ex{2} Razona si son ciertas o falsas las afirmaciones siguientes:\begin{tasks}(1)
  \task Si $\vec u\cdot\vec v=0$ y $\vec u\times\vec v=\vec 0$, entonces alguno de los dos vectores es nulo.
  \task Tres vectores no nulos y perpendiculares dos a dos forman una base.
  \task Si $[\vec u,\vec v,\vec w]=0$, alguno de los tres vectores es nulo.
  \task Dos rectas que no tienen ningún punto en común son paralelas.
  \task Si $\vec u\times\vec v=\vec u\times\vec w$ y $\vec u\neq\vec 0$, entonces $\vec v=\vec w$.
\end{tasks}

\ex{2} Un cubo tiene un vértice en el origen, tres aristas sobre los semiejes positivos y arista 2. Halla la ecuación del plano que pasa por los vértices $A(2,\,0,\,0)$, $B(0,\,2,\,0)$ y $C(0,\,0,\,2)$, el área del triángulo $ABC$ y el punto donde la diagonal del cubo que sale del origen corta este plano.

\ex{2} Un dron sale del punto $P(1,\,2,\,0)$ y vuela en línea recta en la dirección del vector $(2,\,1,\,3)$. En su trayectoria están el tejado inclinado de una nave, contenido en el plano $x+y+z=12$, y un poste vertical que pasa por el punto $(5,\,4,\,0)$. ¿Con cuál de los dos obstáculos choca primero?

\ex{2} \pau{Aragón, extraordinaria 2026} Estamos reparando una silla metálica que cojeaba, porque los extremos inferiores de sus cuatro patas no se apoyaban a la vez en un suelo plano. La silla está tumbada sobre una mesa de trabajo de 120 cm de altura y tiene cuatro patas de 45 cm de longitud cada una. Después de la reparación, los extremos inferiores de las patas están situados en los puntos $A(-2,\,3,\,120)$, $B(20,\,42,\,120)$, $C(-1,\,4,\,165)$ y $D(20,\,43,\,165)$. Cuando bajemos la silla de la mesa y la dejemos en el suelo, ¿cojeará o apoyará las cuatro patas sobre el suelo plano del taller?

\ex{2} \pau{Islas Baleares, extraordinaria 2026} En un estudio de astronomía se modeliza una constelación artificial de satélites que orbitan alrededor de la Tierra. Tomamos un sistema de coordenadas con el origen en el centro de la Tierra (las unidades son decenas de miles de kilómetros) y situamos cinco satélites: $A(1,\,2,\,3)$, $B(3,\,2,\,1)$, $C(5,\,4,\,3)$, $D(2,\,1,\,1)$ y $E(4,\,3,\,1)$. Decimos que unos satélites forman un plano orbital si, junto con el centro de la Tierra, son coplanarios.\begin{tasks}(1)
  \task Determina si los satélites $A$, $B$ y $C$ forman un plano orbital.
  \task Proporciona un punto $F$ (distinto del origen) donde se podría colocar un sexto satélite para que los satélites $D$, $E$ y $F$ formen otro plano orbital.
\end{tasks}

\ex{2} \pau{Islas Baleares, ordinaria 2026} En un estudio de grabación se han colgado del techo tres micrófonos para grabar un concierto acústico. Tomamos un sistema de referencia con el origen en el suelo de la sala; el techo está a 7 m de altura. El micrófono $M_1$ tiene coordenadas $x=2$, $y=3$ y cuelga de un cable de 1 m; el micrófono $M_2$ tiene $x=4$, $y=4$ y cuelga de un cable de 2 m, y el micrófono $M_3$ tiene $x=5$, $y=1$ y cuelga de un cable de 3 m. Determina las coordenadas de los tres micrófonos y calcula la ecuación del plano que los contiene.

\ex{2} \pau{La Rioja, ordinaria 2026} En un chiringuito de playa se instala un toldo para los clientes. El toldo se coloca en los puntos de apoyo $\left(3,\,8,\,\tfrac{17}{2}\right)$, $\left(4,\,6,\,\tfrac{37}{4}\right)$, $(2,\,8,\,8)$ y $(s,\,16,\,8)$, con $s\in\mathbb R$.\begin{tasks}(1)
  \task Determina la ecuación del plano que contiene el toldo.
  \task Calcula el valor de $s$ para que los cuatro puntos sean coplanarios.
\end{tasks}

\vfill
\begin{center}\color{muted}\small\faCheckCircle[regular]\ \ ¡Ánimo! Comprueba tus resultados con las herramientas 3D de los apuntes interactivos.\end{center}

\end{document}
